% DocShield AI — Business Logic
% SUM 72H Build Challenge — Deliverable 04
% Compile: pdflatex 04-business-logic.tex  (twice)

\documentclass[10pt,a4paper]{article}

% tcolorbox charge xcolor sans l'option « table » ; on la passe en amont pour
% que \rowcolor soit disponible.
\PassOptionsToPackage{table}{xcolor}

\usepackage[T1]{fontenc}
\usepackage[utf8]{inputenc}
\usepackage[margin=16mm,top=14mm,bottom=12mm]{geometry}

\usepackage{charter}
\usepackage[activate={true,nocompatibility},final,tracking=true,
            factor=1100,stretch=12,shrink=12]{microtype}

\usepackage{amssymb}
\usepackage{xcolor}
\usepackage{array}
\usepackage{enumitem}
\usepackage[most]{tcolorbox}
\usepackage[hidelinks]{hyperref}

\definecolor{navy}{HTML}{123A5F}
\definecolor{steel}{HTML}{3E6D9C}
\definecolor{good}{HTML}{1B6E4F}
\definecolor{goodbg}{HTML}{F1F8F4}
\definecolor{bad}{HTML}{A8322D}
\definecolor{badbg}{HTML}{FDF3F2}
\definecolor{ink}{HTML}{16202A}
\definecolor{muted}{HTML}{68727E}
\definecolor{zebra}{HTML}{F4F7FA}

\pagestyle{empty}
\setlength{\parindent}{0pt}
\setlength{\parskip}{4pt}
\color{ink}

\setlist[itemize]{leftmargin=13pt,itemsep=2pt,topsep=2pt,parsep=0pt}

%% ------------------------------------------------------ encadrés à onglet
\tcbset{
  callout/.style={
    enhanced, sharp corners=downhill, arc=1.5pt,
    colback=white, boxrule=0.8pt,
    left=6pt, right=6pt, top=8pt, bottom=6pt,
    attach boxed title to top left={xshift=7pt, yshift=-\tcboxedtitleheight/2},
    boxed title style={arc=1pt, boxrule=0pt, left=5pt, right=5pt,
                       top=1.5pt, bottom=1.5pt},
    fonttitle=\footnotesize\bfseries\color{white},
  }
}
\newtcolorbox{note}[1]{callout, colframe=navy, colbacktitle=navy, title={#1}}
\newtcolorbox{goodbox}[1]{callout, colframe=good, colback=goodbg,
  colbacktitle=good, title={#1}}
\newtcolorbox{badbox}[1]{callout, colframe=bad, colback=badbg,
  colbacktitle=bad, title={#1}}

\newcommand{\heading}[2]{%
  \par\vspace{12pt}{\color{navy}\bfseries #1\hspace{6pt}#2}\par\vspace{-2pt}}
\newcommand{\fld}[1]{\textbf{\color{navy}#1}\hspace{5pt}}
\newcommand{\hyp}{{\footnotesize\color{bad}\bfseries[hypothesis]}\hspace{3pt}}

\begin{document}

%% ------------------------------------------------------ en-tête
\noindent
\begin{minipage}[b]{\linewidth}
  \colorbox{navy}{\makebox[7mm][c]{\rule[-2pt]{0pt}{12pt}\color{white}\bfseries 04}}%
  \hspace{5pt}{\Large\bfseries\color{navy}Business Logic}\\[3pt]
  \hspace{12mm}{\color{muted}DocShield AI --- customer, value, model, market, next step}
\end{minipage}

\vspace{5pt}
\textcolor{navy}{\rule{\linewidth}{1pt}}

%% ------------------------------------------------------ 1
\heading{1.}{Customer --- who pays, and who signs}

\fld{User} The document reviewer: a credit analyst, an HR recruiter, a
compliance officer. They feel the pain but rarely hold the budget.

\fld{Buyer} The head of risk, compliance or operations. They do not buy time
savings; they buy \emph{reduced exposure} --- a forged payslip approved, a
diploma never checked, a regulator asking how a file was verified.

\fld{Beachhead} Lending and leasing companies in Morocco. High document volume,
direct financial loss on a miss, and a verification step that already exists ---
so the product replaces a manual task rather than creating a new one.

%% ------------------------------------------------------ 2
\heading{2.}{What two external reviews changed}

The prototype was put in front of two people during the challenge, chosen for
opposite vantage points: an HR practitioner who verifies documents himself, and
a company director who would sign for the tool. Deliverable 03 records the
detail; what follows is only what bears on the business.

\vspace{2pt}
\fld{The user tested the mechanism.} Reading and extraction were judged reliable,
including on a flattened scan. The objection was not that the tool fails to see
--- it was that the score could not yet be trusted. For a screening product, that
is the whole sale.

\fld{The buyer named the criterion.} Unprompted, the director identified the
central challenge as \emph{demonstrating the reliability of the results, and
keeping anomaly detection clearly distinct from proof of authenticity}. He was
describing the product's founding constraint without having been told it.

\vspace{4pt}
\begin{note}{The beachhead was not confirmed --- it was redirected}
Section 1 assumes lending and leasing. The one director asked pointed instead to
\textbf{public administrations, educational institutions and companies}, on his
own initiative.

\vspace{3pt}
One opinion does not move a beachhead: a single data point cannot outweigh the
reasoning behind lending, which rests on direct financial loss per miss. But it
is the only segment signal the challenge produced from someone evaluating the
tool for his own organisation, and it is recorded here rather than discarded for
disagreeing with the hypothesis. It is the first question for the ten interviews
in week 1.
\end{note}

\vspace{4pt}
\fld{What this reorders.} Both reviewers converged on the same purchase
criterion, and it is not detection coverage. It is whether a flag can be
trusted, and whether the product is honest about what it does not prove. That
puts the audit trail and the calibration evidence ahead of new detection
categories on the roadmap --- a ranking neither reviewer was asked about, and
both implied.

%% ------------------------------------------------------ 3
\heading{3.}{Value --- what is actually being bought}

\begin{minipage}[t]{0.485\linewidth}
\begin{goodbox}{Measurable}
Verification time per file. Share of files needing a second reviewer. A written,
timestamped trace of every check performed.
\end{goodbox}
\end{minipage}\hfill
\begin{minipage}[t]{0.485\linewidth}
\begin{badbox}{Not yet measured}
Minutes saved per document, and the annual cost of an undetected forgery. Both
must come from interviews, not from an assumption.
\end{badbox}
\end{minipage}

\vspace{3pt}
The second column is deliberately empty of numbers. Quantifying the value
before speaking to a single reviewer would produce a convincing figure and a
worthless one.

%% ------------------------------------------------------ 4
\heading{4.}{Business model}

\vspace{2pt}
% Les 80mm retranchés couvrent les deux colonnes fixes (71mm) et les quatre
% \tabcolsep des gouttières internes, que la largeur des colonnes ignore.
\begin{tabular}{@{}>{\raggedright\arraybackslash}p{27mm}
                  >{\raggedright\arraybackslash}p{44mm}
                  >{\raggedright\arraybackslash}p{\dimexpr\linewidth-80mm}@{}}
\rowcolor{zebra}\textbf{\color{navy}Tier} & \textbf{\color{navy}For} &
\textbf{\color{navy}Basis} \\[2pt]
Pay per document & Low volume, occasional checks & Per analysis, prepaid credits \\[3pt]
\rowcolor{zebra}Per seat & Teams with a recurring flow & Monthly, per reviewer, fair-use cap \\[3pt]
API \& integration & Case-management systems & Annual contract, volume tiers, SLA \\
\end{tabular}

\vspace{5pt}
\begin{note}{Unit economics --- the one number already known}
Each analysis costs a few cents of model inference, readable directly from the
provider's usage dashboard. Cost scales per document, not per customer, and
falls as models get cheaper. Gross margin is therefore high and predictable;
the real cost centres are sales and the trust-building work, not compute.
\end{note}

%% ------------------------------------------------------ 5
\heading{5.}{Market --- sized bottom-up, not asserted}

Rather than quote a market figure, the opportunity is expressed as a formula
whose inputs are verifiable one by one:

\vspace{3pt}
{\color{navy}\bfseries
Organizations in segment $\times$ reviewers per organization $\times$ price per
seat $\times$ 12}

\vspace{4pt}
\hyp Each input is a research task, not a guess: counts come from professional
registers, reviewer headcount from interviews, price from the first pilot
negotiation. \emph{The formula is filled in only once those three are known.}

\vspace{2pt}
\fld{Expansion path} Lending $\rightarrow$ insurance $\rightarrow$ recruitment
$\rightarrow$ education $\rightarrow$ public administration. Same engine, new
document families and new business rules --- which is also the natural order of
increasing sales-cycle length.

%% ------------------------------------------------------ 6
\heading{6.}{Next 30 days}

\begin{tabular}{@{}>{\raggedright\arraybackslash}p{16mm}
                  >{\raggedright\arraybackslash}p{\dimexpr\linewidth-21mm}@{}}
\rowcolor{zebra}\textbf{\color{navy}Week 1} &
Ten interviews with reviewers and their managers. Measure the current process
before proposing to change it. \\[3pt]
\textbf{\color{navy}Week 2} &
Close the production gaps: rate limiting, payload ceiling, timestamped audit
report. Without them there is no pilot. \\[3pt]
\rowcolor{zebra}\textbf{\color{navy}Week 3} &
Run one organization's real files through the tool, side by side with their
manual check. Compare, and record every disagreement. \\[3pt]
\textbf{\color{navy}Week 4} &
Put a price on the table and see whether anyone pays. Fill in the market
formula with the three inputs now known. \\
\end{tabular}

\vspace{6pt}
\begin{badbox}{What would break this}
If reviewers do not trust a flag they cannot verify themselves, the tool adds a
step instead of removing one. That is the failure mode to test first --- before
detection accuracy, before pricing, before market size.
\end{badbox}


\end{document}
